\documentclass[11pt]{article}
\usepackage[utf8]{inputenc}
\usepackage[T1]{fontenc}
\usepackage{graphicx}
\usepackage{booktabs}
\usepackage{amsmath}
\usepackage{amssymb}
\usepackage[margin=1in]{geometry}
\usepackage{caption}
\usepackage{array}
\begin{document}
\title{FX-217 versus placebo in fictional score syndrome: a synthetic two-arm trial}
\author{Northbridge Clinical Unit and Larkspur Research Ward}
\maketitle
\begin{abstract}
This synthetic report describes a fictional two-arm trial of compound code FX-217.
Investigators screened {{v:n_screened}} adults and randomized {{v:n_randomized}}
participants to FX-217 (n={{v:n_alloc_fx}}) or placebo (n={{v:n_alloc_pb}}).
The analyzed sets contained {{v:n_analyzed_fx}} and {{v:n_analyzed_pb}} participants.
At week 8 the mean Fictional Symptom Score change was {{v:chg_w8_fx}} with FX-217 and
{{v:chg_w8_pb}} with placebo (difference {{v:diff_w8}}, 95\% CI {{v:diff_w8_lo}} to
{{v:diff_w8_hi}}, p={{v:p_primary}}).
\end{abstract}
\section{Introduction}
Fictional score syndrome is an invented symptom construct used only so this example can carry a complete randomized-trial skeleton.
Published discussions of made-up symptom scales describe how a 0--100 score can be summarized without referring to a real disease \cite{hale2018,okonkwo2020}.
Placebo change on the Fictional Symptom Score has been sketched in the same imaginary literature \cite{brant2021}.
Flow diagrams, proportion wording, and interval reporting for synthetic studies are taken from those fictional notes \cite{voss2019,chen2022,ibrahim2017}.
FX-217 is a compound code created for this case. It is not a marketed product, and the sites named below are not real hospitals.
The example keeps one clean numeric source so a later typesetting step can emit a clean PDF and a defect PDF that differ only where an override says they differ.
No sponsor, registry entry, or external identifier is attached to the case.
\section{Methods}
The study is synthetic. No real participants were enrolled.
Two imaginary sites took part: Northbridge Clinical Unit and Larkspur Research Ward.
Adults described as having fictional score syndrome were screened ({{v:n_screened}}).
Those who remained after screening were randomized ({{v:n_randomized}}) 1:1 to FX-217 ({{v:n_alloc_fx}}) or placebo ({{v:n_alloc_pb}}) by a fictional allocation list.
The Fictional Symptom Score runs from 0 to 100. A negative change means improvement on that invented scale.
The primary comparison was the week-8 mean change, reported as a difference with a 95\% CI and a p value, without a test statistic.
A response was a decrease of at least ten points. Adverse events were counted once per analyzed participant.
The write-up uses only the values bound in the case. It does not query an external index.
\section{Results}
Of {{v:n_screened}} people screened, {{v:n_excluded}} were excluded
({{v:n_excluded_ineligible}} did not meet criteria and {{v:n_excluded_declined}} declined).
{{v:n_randomized}} participants were randomized, {{v:n_alloc_fx}} to FX-217 and {{v:n_alloc_pb}} to placebo.
Discontinuation occurred in {{v:n_disc_fx}} FX-217 participants
({{v:n_disc_fx_ae}} after an adverse event and {{v:n_disc_fx_wd}} by withdrawal)
and in {{v:n_disc_pb}} placebo participants
({{v:n_disc_pb_ae}} after an adverse event and {{v:n_disc_pb_wd}} by withdrawal).
The analyzed sets were {{v:n_analyzed_fx}} and {{v:n_analyzed_pb}}.
Baseline characteristics follow.
\falsifytable{T1}
Week-4 mean changes were {{v:chg_w4_fx}} with FX-217 and {{v:chg_w4_pb}} with placebo.
Week-8 mean change was {{v:chg_w8_fx}} (95\% CI {{v:chg_w8_fx_lo}} to {{v:chg_w8_fx_hi}}) with FX-217
and {{v:chg_w8_pb}} (95\% CI {{v:chg_w8_pb_lo}} to {{v:chg_w8_pb_hi}}) with placebo.
The difference was {{v:diff_w8}} (95\% CI {{v:diff_w8_lo}} to {{v:diff_w8_hi}}, p={{v:p_primary}}).
Response was recorded for {{v:resp_fx_k}} of {{v:n_analyzed_fx}} FX-217 participants ({{v:resp_fx_pct}}\%)
and {{v:resp_pb_k}} of {{v:n_analyzed_pb}} placebo participants ({{v:resp_pb_pct}}\%),
a difference of {{v:resp_diff}} percentage points.
Any adverse event was recorded for {{v:ae_fx_k}} of {{v:n_analyzed_fx}} ({{v:ae_fx_pct}}\%)
and {{v:ae_pb_k}} of {{v:n_analyzed_pb}} ({{v:ae_pb_pct}}\%).
\falsifytable{T2}
\falsifyfigure{F1}
\section{Discussion}
{{v:ph_concl_claim}}
{{v:ph_hedge}}
{{v:ph_limitation}}
The flow counts add across exclusion, allocation, discontinuation, and analysis, and each printed percentage is the corresponding count divided by its denominator.
Confidence intervals contain their point estimates, and the week-8 difference equals the FX-217 change minus the placebo change.
Those relations are properties of this fictional arithmetic, not findings about a treatment.
Readers should not treat FX-217, Northbridge Clinical Unit, or Larkspur Research Ward as real \cite{chen2022}.
A later override can change one surface and leave the others untouched, which is why several headline values are repeated in the abstract, the results, and a table or figure.
\falsifybibliography
\end{document}
